Distributed satellite systems can shorten the response to time-critical Earth-observation events, but their performance depends on orbital access, finite communication opportunities, and the routing policy used to disseminate tasks. This thesis studies how assigning a fraction of a Walker constellation to a central-node communication and routing role affects wildfire-task dissemination, follow-up observation, resource demand, and failure sensitivity. A traceable preprocessing pipeline converts de-duplicated Sentinel-3 SLSTR fire-radiative-power detections into 805 prioritized tasks: 32 for California and 773 for India. Each task retains a release time, location, priority-dependent data volume, deadline, and source provenance.

The architecture grid combines 60, 120, and 180 satellites; 3, 5, and 10 planes; altitudes of 500, 800, and 1000 km; inclinations of 60, 80, and 98.6 degrees; and central-node fractions from 0 to 100 percent. The earlier nominal campaign contains 891 validated cases per region and establishes that architecture preferences are region- and metric-dependent. The newer Full891 Fresh V2 campaign evaluates 539 operational, policy, task-size, and failure scenarios per architecture. The locally available California archive contains 890 complete cases and 480,110 unique case-scenario rows after resume de-duplication; all 19,179 architecture checks pass.

For California Fresh V2, the bounded nominal policy achieves mean observation fulfilment of 61.59 percent, strict useful-recipient dissemination of 0.63 percent, and useful-recipient coverage of 15.29 percent. Removing hop and copy restrictions in the \texttt{unlimited\_useful\_deadline} diagnostic increases these means to 77.62, 74.69, and 91.14 percent, at the cost of increasing transferred data from 197.55 to 2,047.57 Mbit per case. Under the nominal policy, observation fulfilment is non-monotonic: it reaches 71.13 percent at 50 percent central nodes in the balanced subset and falls to 35.70 percent at 100 percent. Under the unlimited diagnostic, dissemination saturates while communication burden continues to rise.

At a 30 percent stress level, ground-station outage is the dominant California vulnerability, reducing observation and strict dissemination by 21.46 and 20.84 percentage points on average. Random satellite and central-node failures produce smaller mean losses, while capacity degradation has limited average effect. Increasing task size from one to four times the nominal value reduces mean observation fulfilment from 77.62 to 69.48 percent and strict dissemination from 74.69 to 60.10 percent. A Pareto-consistent objective-weight analysis selects a 60-satellite, 10-plane, 800 km, 98.6-degree, 40-percent-central-node configuration for balanced and performance-dominant preferences; it wins 41.79 percent of 10,000 randomized preference draws. This is a conditional compromise, not a universal optimum. India Fresh V2 and observation-radius sensitivity remain required before the final cross-region robustness conclusion.
